\chapter{1913:Alfred Werner}
\label{chap:chemistry1913}

The Nobel Prize in Chemistry for the year 1913 was awarded to Alfred Werner.

\textbf{Motivation:}
in recognition of his work on the linkage of atoms in molecules by which he has thrown new light on earlier investigations and opened up new fields of research especially in inorganic chemistry

\section{Alfred Werner}

\subsection*{Early Life and Education}

Alfred Werner was born on 1866-12-12 in Mulhouse, France.
Alfred Werner was a Swiss chemist who was a student at ETH Zurich and a professor at the University of Zurich. He won the Nobel Prize in Chemistry in 1913 for proposing the octahedral configuration of transition metal complexes. Werner developed the basis for modern coordination chemistry. He was the first inorganic chemist to win the Nobel Prize, and the only one prior to 1973.

\subsection*{Career and Major Research}

Werner was professor of chemistry at the University of Zürich. In 1893 he proposed the coordination theory, according to which metal atoms exhibit a primary valence and a secondary valence that fixes the spatial arrangement of surrounding groups. This theory explained the structure of complex salts and the existence of optical isomers in coordination compounds.

\subsection*{Nobel Prize-winning Work}

The work for which Alfred Werner was awarded the Nobel Prize in Chemistry in 1913 was described by the Nobel Committee as: "in recognition of his work on the linkage of atoms in molecules by which he has thrown new light on earlier investigations and opened up new fields of research especially in inorganic chemistry".

This research fundamentally advanced the field by making groundbreaking contributions that transformed chemical science.

\subsection*{Significance and Impact}

The impact of Alfred Werner's work extends far beyond the specific achievement recognized by the Nobel Prize.
These contributions continue to influence chemical research and education today.

\subsection*{Later Career and Honors}

Alfred Werner continued his scientific work until 1919-11-15, passing away in Zurich, Switzerland.
He received numerous honors including membership in major scientific academies and societies.

\subsection*{Legacy}

The legacy of Alfred Werner endures through the lasting impact of his scientific contributions.
Alfred Werner's work continues to be cited and built upon by researchers across the chemical sciences.

\subsection*{Modern Developments}

Werner's coordination theory explains the structure of coordination compounds and remains the framework for modern inorganic chemistry, bioinorganic chemistry, and catalysis. Werner-type complexes are central to organometallic chemistry, homogeneous catalysis (including the Nobel-winning metathesis and cross-coupling reactions), and the design of MRI contrast agents and anticancer drugs such as cisplatin.

\subsection*{Selected Publications}

Alfred Werner's major publications include numerous papers in leading journals of the era, detailing the experimental and theoretical work summarized above.

\clearpage

\section*{Chapter Summary}

Alfred Werner was awarded the prize for his work on the linkage of atoms in molecules, which established the field of coordination chemistry.
